% --- Chapter appendices: numbered A.A, A.B ----------------------------------
\renewcommand{\thesection}{\thechapter.\Alph{section}}
\setcounter{section}{0}

\section{Elementary Koenigs functions at \texorpdfstring{$r=2$}{r=2} and \texorpdfstring{$r=4$}{r=4}}
\label{a:app:integrable}

These cases lie outside subcritical branching ($r=2p<1$) but complete the
classical picture of integrable logistic linearisations. Both are recorded in full
because they show concretely what a solvable case looks like, and because they sit
exactly where the Becker--Bergweiler classification, \cref{a:thm:BB}, says all
differentially algebraic Schr\"oder functions must sit: over repelling fixed
points.

\subsection{Case \texorpdfstring{$r=2$}{r=2} (logarithmic)}
\label{a:app:r2}

For $r=2$, $f_2(x)=2x(1-x)$. The fixed point $0$ has multiplier $\lambda=2$
(repelling); the linearisation is of Poincar\'e type rather than an
attracting-Koenigs function, but an elementary closed form still exists.

The affine change $u=1-2x$ conjugates the dynamics to $u\mapsto u^2$, since
\begin{equation*}
u_{n+1} = 1-2\bigl(2x_n(1-x_n)\bigr) = 1-4x_n+4x_n^2 = (1-2x_n)^2 = u_n^2 .
\end{equation*}
Linearising $u\mapsto u^2$ by the logarithm, since $\ln(u^2)=2\ln u$, and
returning to $x$ gives the candidate
\begin{equation}
\label{a:eq:psi2}
\psi_2(x) = -\tfrac12\ln(1-2x),
\end{equation}
normalised so that $\psi_2(0)=0$ and $\psi_2'(0)=1$. Indeed,
\begin{equation*}
\psi_2\bigl(f_2(x)\bigr)
= -\tfrac12\ln\bigl(1-4x+4x^2\bigr)
= -\tfrac12\ln\bigl((1-2x)^2\bigr)
= 2\,\psi_2(x) .
\end{equation*}

\subsection{Case \texorpdfstring{$r=4$}{r=4} (inverse trigonometric)}
\label{a:app:r4}

For $r=4$, $f_4(x)=4x(1-x)$. The fixed point $0$ has multiplier $f_4'(0)=4$
(repelling); as with $r=2$, this is a Poincar\'e-type linearisation, not an
attracting Koenigs function. With $x=\sin^2\theta$ one finds
\begin{equation*}
f_4(x) = 4\sin^2\theta\cos^2\theta = (2\sin\theta\cos\theta)^2 = \sin^2(2\theta),
\end{equation*}
so the map is a doubling $\theta\mapsto2\theta$ in the $\theta$ coordinate. The
candidate
\begin{equation}
\label{a:eq:psi4}
\psi_4(x) = \bigl(\arcsin\sqrt{x}\bigr)^2
\end{equation}
satisfies $\psi_4(f_4(x))=4\psi_4(x)$ on a suitable real domain: using
$\arcsin\bigl(2u\sqrt{1-u^2}\bigr)=2\arcsin u$ with $u=\sqrt{x}$,
\begin{equation*}
\psi_4\bigl(f_4(x)\bigr)
= \bigl(\arcsin\sqrt{4x(1-x)}\bigr)^2
= \bigl(2\arcsin\sqrt{x}\bigr)^2
= 4\,\psi_4(x) .
\end{equation*}
And as $x\to0$, $\arcsin\sqrt{x}\approx\sqrt{x}$, so $\psi_4(x)\approx x$ and
$\psi_4'(0)=1$.

\subsection{Location inside the classification}
\label{a:app:location}

Under the affine conjugacy of \cref{a:lem:affine}, $r=2\Leftrightarrow c=0$
(monomial $z^2$) and $r=4\Leftrightarrow c=-2$ (Chebyshev type; the value $c=-2$
is also reached at $r=-2$). Neither exceptional value lies in the image of
$r\in(0,1)$, which is $c\in(0,\tfrac14)$. As explained in \cref{a:rem:empty},
this parameter count is a consistency check rather than the proof: the decisive
point is that both exceptional cases sit over \emph{repelling} fixed points, which
the attracting window excludes structurally.

The two derivations above realise the first two families of \cref{a:thm:BB}
exactly. For $r=2$, inverting \cref{a:eq:psi2}: $t=-\tfrac12\ln(1-2x)$ gives
\begin{equation*}
\varphi_2(t) = \psi_2^{-1}(t) = \frac{1-\ee^{-2t}}{2},
\end{equation*}
an affine (hence M\"obius) image of $\ee^{-2t}$ --- the family
$\exp(\alpha t^{\rho})$ with $\rho=1$, $\alpha=-2$. Correspondingly the
substitution $u=1-2x$ conjugates $f_2$ to the monomial $u\mapsto u^2$, and the
linearised fixed point has multiplier $2$ (repelling), as \cref{a:thm:BB}(i)
requires. For $r=4$, inverting \cref{a:eq:psi4}: $t=(\arcsin\sqrt{x})^2$ gives
\begin{equation*}
\varphi_4(t) = \psi_4^{-1}(t) = \sin^2\!\sqrt{t}
= \frac{1-\cos\bigl(2\sqrt{t}\bigr)}{2},
\end{equation*}
a M\"obius image of $\cos(\alpha t^{\rho}+\beta)$ with $\rho=\tfrac12$,
$\alpha=2$, $\beta=0$ --- the fractional exponent is admissible precisely because
$\cos(2\sqrt{t})$ is an entire function of $t$. Correspondingly $u=1-2x$ conjugates
$f_4$ to $T_2(u)=2u^2-1$, the second Chebyshev polynomial, and the multiplier at
the linearised fixed point is $4$ (repelling).

For $r\in(0,1)$ the fixed point is attracting, no member of the classification is
available, and \cref{a:thm:mainht} gives hypertranscendence of $\psi_r$ --- which
underpins the claims of this thesis for $A(p)=2\Psi_{2p}(\tfrac12)$.
